% Dados de identificação do trabalho (capa, folha de rosto, metadados do PDF).
% Equivale ao 00-dados.tex do modelo ufpr-abntex.

\tipotrabalho{Trabalho de Conclusão de Curso}

\titulo{Impactos da Inteligência Artificial no Mercado de Trabalho:
        uma Revisão Sistemática da Literatura (2013--2026)}
\autor{Acácio Vasconcelos Telechi}                      % TODO: nome completo
\local{Curitiba}
\data{2026}                         % apenas o ano

% Preencher \orientador OU \orientadora (se ambos, prevalece o primeiro).
\orientador{Victor Rodrigues de Oliveira}   % TODO: ex.: Prof. Dr. Fulano de Tal
\orientadora{}
\coorientador{}
\coorientadora{}
\scoorientador{}
\scoorientadora{}

\instituicao{Universidade Federal do Paraná}
\def\ImprimirSetor{Setor de Ciências Sociais Aplicadas}
\def\ImprimirProgramaPos{}
\def\ImprimirCurso{Curso de Ciências Econômicas}

% O ponto final é acrescentado pelo modelo.
\preambulo{Trabalho de Conclusão de Curso apresentado ao curso de Graduação em
Ciências Econômicas, Setor de Ciências Sociais Aplicadas, Universidade Federal
do Paraná, como requisito parcial à obtenção do título de Bacharel em Ciências
Econômicas}

\hypersetup{
    pdftitle={Impactos da Inteligência Artificial no Mercado de Trabalho: uma Revisão Sistemática da Literatura (2013-2026)},
    pdfauthor={\imprimirautor},
    pdfsubject={\imprimirtipotrabalho},
    pdfkeywords={inteligência artificial; IA generativa; mercado de trabalho; emprego; abordagem de tarefas; revisão sistemática},
    urlcolor=blue!50!black,
}
